\section*{Capítulo \ref{ch:b1:electrophilic-additions} --- Alcenos e alcinos: adições eletrofílicas}

\begin{solution}{exo:b1:electrophilic-additions:1}
2-Cloropropano; 2-bromo-2-metilpropano; 1-bromo-1-metilciclo-hexano;
2-clorobut-2-eno.
\end{solution}

\begin{solution}{exo:b1:electrophilic-additions:2}
1,2-Dibromoetano; \textit{trans}-1,2-dibromociclo-hexano;
1,2-dibromopropano.
\end{solution}

\begin{solution}{exo:b1:electrophilic-additions:3}
Propan-2-ol; 2-metilpropan-2-ol; 1-metilciclo-hexan-1-ol (Markovnikov em
cada caso).
\end{solution}

\begin{solution}{exo:b1:electrophilic-additions:4}
Propino (\omterm{def:b1:electrophilic-additions:acetylide}{alcino terminal}): \ce{CH3C#CH + NH2- -> CH3C#C- + NH3}. Etanol:
\ce{CH3CH2OH + NH2- -> CH3CH2O- + NH3}. O but-2-ino e o propeno não têm
hidrogênio ácido o bastante.
\end{solution}

\begin{solution}{exo:b1:electrophilic-additions:5}
O par $\pi$ capta \ce{H+} no carbono \ce{CH2}: cátion terciário
\ce{(CH3)3C+}; um \omterm{def:b1:lewis-resonance-vsepr:lewis}{par não ligante} da água se liga a ele: \ce{(CH3)3COH2+};
uma molécula de água retira um próton: \ce{(CH3)3COH}. Em ácido
concentrado quente, a água é escassa e o alceno volátil escapa: a reação
inversa, a \omterm{def:b1:alcohol-activation:dehydration}{desidratação}, avança.
\end{solution}

\begin{solution}{exo:b1:electrophilic-additions:6}
A protonação do 2-metilpropeno dá um cátion terciário, a do propeno um
secundário. A protonação determinante tem um \omterm{def:b1:elementary-steps:energy-profile}{estado de transição}
parecido com o cátion (Hammond): o cátion terciário, mais estável, é
alcançado por uma barreira mais baixa.
\end{solution}

\begin{solution}{exo:b1:electrophilic-additions:7}
O íon bromônio faz uma ponte sobre uma face do anel; o brometo o abre
pela outra face: os dois bromos terminam em \textit{trans}. O dibrometo
\textit{trans} é \omterm{def:b1:stereochemistry-in-depth:stereocentre}{quiral} (sem plano de simetria), mas o brometo ataca os
dois carbonos igualmente: os dois \omterm{def:b1:stereochemistry-in-depth:enantiomers}{enantiômeros} se formam em quantidades
iguais, e o produto é racêmico, opticamente inativo.
\end{solution}

\begin{solution}{exo:b1:electrophilic-additions:8}
A protonação de C1 dá o cátion secundário em C2, vizinho do C3 terciário
que leva um hidrogênio; esse hidrogênio migra com seu par para C2 (um
deslocamento 1,2 de hidreto), deixando um cátion terciário em C3,
capturado por \ce{Cl-}: 2-cloro-2-metilbutano.
\end{solution}

\begin{solution}{exo:b1:electrophilic-additions:9}
Propino + amideto de sódio: o íon propineto; com o 1-bromopropano (SN2):
\ce{CH3C#C-CH2CH2CH3}, hex-2-ino. O propineto com o benzaldeído, depois
ácido diluído: \ce{C6H5CH(OH)C#CCH3}, 1-fenilbut-2-in-1-ol.
\end{solution}

\begin{solution}{exo:b1:electrophilic-additions:10}
\cip{E}-but-2-eno: bromônio em uma face, brometo pela outra em C2 ou em
C3; os dois ataques dão o mesmo composto, de descritores \cip{2R,3S}: o
\omterm{def:b1:stereochemistry-in-depth:enantiomers}{composto meso}, um único \omterm{def:b1:stereochemistry-in-depth:isomers}{estereoisômero}. \cip{Z}-but-2-eno: o ataque em C2
dá \cip{2S,3S} (ou sua imagem especular a partir da outra face), o ataque
em C3 \cip{2R,3R}; quantidades iguais: dois \omterm{def:b1:stereochemistry-in-depth:isomers}{estereoisômeros}, uma \omterm{def:b1:stereochemistry-in-depth:enantiomers}{mistura racêmica}.
\end{solution}

\begin{solution}{exo:b1:electrophilic-additions:11}
No íon bromônio assimétrico, o carbono mais substituído leva mais da
carga positiva (sua ligação C--Br é mais longa e mais fraca, como em um
cátion quase terciário ou secundário). A água, em grande excesso, ataca
esse carbono pela face oposta à ponte: 1-bromopropan-2-ol.
\end{solution}

\begin{solution}{exo:b1:electrophilic-additions:12}
$D = 1$: um alceno. O 2-metilbut-1-eno e o 2-metilbut-2-eno dão ambos o
cátion terciário, logo o 2-bromo-2-metilbutano e o 2-metilbutan-2-ol. Seus
espectros de RMN de próton diferem: dois prótons vinílicos (\ce{=CH2},
singletos) para o primeiro, um (\ce{=CH-}, um quarteto) para o segundo.
\end{solution}

\begin{solution}{pb:b1:electrophilic-additions:1}
\textbf{1.} Em cada carbono, \ce{CH3} tem prioridade sobre H: o \cip{Z}
tem os dois metilas do mesmo lado, o \cip{E} em lados opostos.
\textbf{2.} \omterm{def:b1:stereochemistry-in-depth:isomers}{Estereoisômeros} que não são imagens especulares: \omterm{def:b1:stereochemistry-in-depth:enantiomers}{diastereoisômeros}.
\textbf{3.} A rotação em torno de C=C romperia a ligação $\pi$, o que custa
muito mais energia do que as colisões fornecem à temperatura ambiente.
\textbf{4.} O \cip{E}: seus grupos metila não se congestionam.
\textbf{5.} \ce{C4H8 + Br2 -> C4H8Br2}.
\textbf{6.} Dois estereocentros equivalentes: no máximo quatro, de fato
três (\cip{2R,3R}, \cip{2S,3S} e o meso \cip{2R,3S}).
\textbf{7.} O par $\pi$ ataca um átomo de Br de \ce{Br2}, o par Br--Br sai
como \ce{Br-}, e um \omterm{def:b1:lewis-resonance-vsepr:lewis}{par não ligante} do bromo atacado se liga ao segundo
carbono: um íon bromônio de três átomos.
\textbf{8.} No íon em ponte, todo átomo tem um octeto; um \omterm{def:b1:electronic-effects:intermediates}{carbocátion}
secundário aberto teria um carbono com seis elétrons.
\textbf{9.} Um \omterm{def:b1:lewis-resonance-vsepr:lewis}{par não ligante} de \ce{Br-} se liga a um carbono do anel; a
ligação C--Br da ponte se rompe, e seu par vai para o bromo da ponte.
\textbf{10.} A ponte ocupa uma face; como na SN2, o \omterm{def:b1:electronic-effects:nucleophile}{nucleófilo} entra do
lado oposto à ligação que se rompe.
\textbf{11.} Uma \omterm{def:b1:electrophilic-additions:halonium}{adição anti}.
\textbf{12.} A água competiria com o brometo pelo íon bromônio e daria uma
bromo-hidrina.
\textbf{13.} \cip{2R,3S} (ou \cip{2S,3R}, o mesmo composto).
\textbf{14.} O mesmo composto.
\textbf{15.} Em uma \omterm{def:b1:stereochemistry-in-depth:isomers}{conformação} com os dois bromos eclipsados, um plano de
simetria passa entre C2 e C3; os descritores são opostos.
\textbf{16.} \cip{2R,3R} e \cip{2S,3S}.
\textbf{17.} Quantidades iguais: uma \omterm{def:b1:stereochemistry-in-depth:enantiomers}{mistura racêmica}.
\textbf{18.} Sim: os dois alcenos diastereoisoméricos dão produtos
diferentes (meso contra racêmico).
\textbf{19.} Nenhuma: os dois \omterm{def:b1:stereochemistry-in-depth:enantiomers}{enantiômeros} se compensam.
\textbf{20.} $5.60/56.0 = \qty{0.100}{mol}$ de cada.
\textbf{21.} $0.100 \times 0.85 \times 215.8 = \qty{18.3}{g}$ de cada.
\textbf{22.} \textbf{$[\alpha] = 0^\circ$ (\omterm{def:b1:stereochemistry-in-depth:enantiomers}{composto meso}), \qty{18.3}{g}}.
\end{solution}
